\section{PONER}

The Python part of this exercise is available in \myhref{https://github.com/LosDelDGIIM/LosDelDGIIM.github.io/blob/main/subjects/Erasmus-DUE/GKI/Exercises/Exercise2.ipynb}{this Jupyter Notebook}.

% // TODO: Corregir

\subsubsection*{Kahn’s Algorithm (Topological Sorting)}

You are responsible for the course planning at a university. For a specific module, there are several lectures that partially build upon each other. Your task is to find an order in which all lectures can be attended without attending a lecture before one of its prerequisites.
\begin{description}
    \item[Lectures (Nodes):] Math 1, Math 2, Computer Science 1, Computer Science 2, Databases, Algorithms, Software Development
    \item[Dependencies (directed edges ``must be attended before''):]
    \begin{itemize}
        \item Math 1 $\rightarrow$ Math 2
        \item Computer Science 1 $\rightarrow$ Computer Science 2
        \item Computer Science 1 $\rightarrow$ Algorithms
        \item Math 2 $\rightarrow$ Algorithms
        \item Computer Science 2 $\rightarrow$ Databases
        \item Algorithms $\rightarrow$ Software Development
    \end{itemize}
\end{description}

\begin{ejercicio}
    Create the graph. Represent the directed graph as an adjacency list, and for each node, note the in-degree (number of incoming edges).
\end{ejercicio}

\begin{ejercicio}
    Apply Kahn's Algorithm. Perform Kahn’s algorithm manually (or implement it in a programming language of your choice) to find a topological sorting of the lectures. Document each step:
    \begin{enumerate}
        \item Start with all nodes whose in-degree $= 0$.
        \item Remove one of these nodes, add it to the order, and reduce the in-degrees of its neighbors.
        \item Repeat until all nodes are processed.
    \end{enumerate}
\end{ejercicio}

\begin{ejercicio}
    Interpret the result. State the calculated order of the lectures, and determine if this order is unique. Justify your answer.
\end{ejercicio}